\documentclass[12pt]{article}
\usepackage[utf8]{inputenc}
\usepackage[T1]{fontenc}
\usepackage{hyperref}
\title{Hausa Sign Language Language Resources for K-12 Curriculum Development}
\author{Compiled by Claude Sonnet 5.5 (claude-sonnet-5-5) for Andr\'as Kornai}
\date{}
\begin{document}
\maketitle

\section{Introduction}
Hausa Sign Language (ISO 639-3 hsl) is the indigenous sign language of Deaf communities in northern Nigeria. Web searches found only a small number of resources, and no K-12 curriculum in the language. Estimates of the number of signers vary widely, and the language is reported not to be formally taught in schools; children usually acquire it from parents and peers.

\section{Primary Educational Resources}
\begin{itemize}
\item Kano-based thematic picture dictionaries of Hausa Sign Language: a project working with Deaf people in Kano State since 2011 reports nine published sign books, distributed free to schools for deaf children, adult literacy classes and deaf associations. \url{https://hausa-sign-language.com/hausa-sign-language-2/}
\item \emph{Maganar hannu / Hausa Sign Language}, a booklet series from Helmut Buske Verlag pairing drawings of signs with Hausa glosses and English translations. \url{https://buske.de/hbv_en/series/Maga}
\end{itemize}

\section{Linguistic and Reference Materials}
\begin{itemize}
\item Wikipedia overview of Hausa Sign Language. \url{https://en.wikipedia.org/wiki/Hausa_Sign_Language}
\item ``The Incorporation of Nigerian Signs in Deaf Education in Nigeria'' (Peter Lang), relevant to Nigerian sign use in deaf schooling; it is not specific to Hausa Sign Language. \url{https://www.peterlang.com/document/1095191}
\end{itemize}

\section{Religious and Community Resources}
Sani Muhammad Jimada, a Deaf advocate in Kano, is described as leading a project to build a Hausa Sign Language dictionary. \url{https://wikitongues.org/fellows/sani-muhammad-jimada/}

\section{Digital and Audio Resources}
The project site above hosts the sign-book material. The search results describe a planned digital dictionary with video, but I did not verify that it is available.

\section{Curriculum Development Context}
Little exists. A K-12 curriculum would have to be built largely from scratch, starting from the Kano picture dictionaries and a video dictionary if one is completed, with Deaf community leadership.

\end{document}